\chapter{Flame Stretch, Curvature and Intrinsic Instabilities}\label{ch:stretch}

The planar, unstretched flame of Chapter~\ref{ch:premixed} is an idealisation. Real flames are curved and sit in non-uniform flow, and a smooth flame may spontaneously wrinkle itself even in a quiescent mixture. Both effects are captured by the concept of \textbf{stretch} and by the theory of intrinsic flame instabilities, which together explain how the local burning velocity deviates from $\sLz$, how flames extinguish, and why hydrogen flames behave so differently. The treatment follows Lecture~9 and its recording, the stretched-flame review of Law, and the instability review of Oppong \emph{et al.}~\cite{lecture9,transcript_l9,Law1988,oppong2022instability}.

\section{Definition of stretch}

Flame stretch $K$ (or $\kappa$) is the fractional rate of change of an element of flame surface area,
\begin{equation}
K=\frac{1}{A}\frac{\dd A}{\dd t}=\underbrace{\grad_t\!\cdot\vv w_t}_{\text{tangential strain}}+\underbrace{(\vv w\cdot\vv n)(\grad\!\cdot\vv n)}_{\text{curvature + motion}},
\label{eq:stretch}
\end{equation}
the sum of two contributions: a \emph{tangential strain} due to flow non-uniformities along the surface, and a term combining \emph{curvature} with the flame's motion (for a steady flame only the curvature part survives)~\cite{lecture9,Law1988}. Velocity gradients \emph{normal} to the flame are excluded from the definition of stretch, although, as we shall see, they still affect the flame by shifting its position~\cite{transcript_l9}. Stretch can be positive (area increasing, e.g.\ an expanding spherical flame or a flame strained in counterflow) or negative (area decreasing), with the sign set by the geometry; the canonical examples are the simple strain $K=\tfrac12(\dd v/\dd x+\dd u/\dd y)$ and the curved-flame value $K=-u\sin\theta/(2R_f)$~\cite{lecture9}. Because a real flame has finite thickness, an iso-surface must be chosen to evaluate $K$~\cite{lecture9}. The two configurations used to study stretch experimentally are the counterflow (opposed-jet) flame, usually fed reactants against products to help sustain it, and the spherically expanding flame~\cite{transcript_l9}.

\section{Stretch on an infinitely thin flame: the hydrodynamic limit}

If the flame is treated as infinitely thin, stretch can act only through the change in surface area, i.e.\ through tangential strain~\cite{lecture9,transcript_l9}. Increasing the area brings more reactants into contact with the flame, so more fuel is burned per unit time, but the local \emph{mass} burning rate per unit area is unchanged --- only the \emph{volumetric} burning rate increases. This is \textbf{hydrodynamic stretch}, a purely geometric effect. Non-dimensionalising the stretch by the chemical time $\tau_{chem}=\dL/\sLz$ gives the \textbf{Karlovitz number},
\begin{equation}
\Ka=\frac{K}{1/\tau_{chem}}=\frac{\dL}{\sLz}\,K ,
\label{eq:karlovitz-stretch}
\end{equation}
so the Karlovitz number is precisely a normalised stretch rate~\cite{lecture9,transcript_l9}. Increasing turbulence intensity increases the stretch and hence $\Ka$, which is the link back to the regime diagram of Chapter~\ref{ch:turbulence}.

\section{Finite-thickness flames: moderate versus large stretch}

A flame of finite thickness responds to stretch in two regimes~\cite{lecture9,transcript_l9}. Under \emph{moderate} stretch the tangential strain perturbs the preheat zone only slightly, and the normal velocity gradient mainly shifts the flame to a new position where the burning velocity matches the local flow; if the flame can reposition itself it maintains complete combustion, so its thickness, temperature and burning velocity stay nearly constant. Whether a change in strain changes the residence time --- measured by the Damk\"ohler number $\Da=\tau_{flow}/\tau_{chem}$ --- thus depends on whether the flame is free to move~\cite{lecture9}. Under \emph{large} stretch, when the flame is restrained (for example pressed against the stagnation plane in a counterflow), the reaction zone itself is affected: the strain-induced convective fluxes carry heat and radicals out of the reaction zone faster than the chemistry can replenish them, and the flame extinguishes through incomplete combustion~\cite{lecture9,transcript_l9}. Typical extinction stretch rates are of order \SI{500}{s^{-1}} for hydrocarbons but up to \SI{5000}{s^{-1}} for hydrogen, which is far more resistant to strain-induced extinction~\cite{transcript_l9}.

\begin{keybox}
\keyhead{When does stretch matter?}
Stretch affects a premixed flame only if the flame is \emph{restrained} from repositioning, or if there is \emph{differential diffusion}; otherwise the flame simply ``runs away'' from the stretch and its structure is unchanged~\cite{lecture9}.
\end{keybox}

\section{Markstein length and differential diffusion}

In the low-stretch limit the stretched burning velocity varies linearly with stretch,
\begin{equation}
\sLstr=\sLz-\mathcal{L}\,K ,
\label{eq:markstein}
\end{equation}
where the \textbf{Markstein length} $\mathcal{L}$, of order the flame thickness $\delta_L$, measures the sensitivity of the flame to stretch and encodes the effect of heat release and curvature on the flame topology~\cite{lecture9,Law1988}. Its sign is controlled by the Lewis number of the deficient reactant, i.e.\ by differential (preferential) diffusion~\cite{transcript_l9,lecture9}:
\begin{itemize}[leftmargin=1.5em,itemsep=2pt]
\item With \emph{equal diffusivities} ($\Le=1$ for all species), the burning velocity and peak temperature stay essentially constant as stretch increases, dropping sharply only at the critical extinction stretch.
\item With $\Le<1$ --- the deficient reactant diffuses into the reaction zone faster than heat is conducted away, as for lean hydrogen ($\phi\approx0.5$--$0.8$) --- the burning velocity and peak temperature \emph{increase} with stretch, producing super-adiabatic temperatures and a \emph{negative} Markstein length, before eventually falling toward extinction.
\item With $\Le\gtrsim1$ --- as for most other cases, including rich mixtures where the deficient reactant is oxygen --- the burning velocity \emph{decreases} monotonically with stretch (positive Markstein length).
\end{itemize}
The Markstein relation~\eqref{eq:markstein} is the linear low-stretch asymptote; the real $\sLstr(K)$ curve deviates from it at higher stretch~\cite{lecture9,transcript_l9}. The strong negative Markstein behaviour of lean hydrogen is the root cause of its tendency to wrinkle and self-accelerate, treated next.

\begin{figure}[htbp]
\centering
\includegraphics[width=0.6\linewidth]{fig_stretch_markstein}
\caption{Stretched laminar flame speed versus stretch rate, illustrating the Markstein relation $\sLstr=\sLz-\mathcal{L}K$ and its dependence on the Lewis number~\cite{lecture9}.}
\label{fig:fig_stretch_markstein}
\end{figure}

\section{Stretch in diffusion flames}

For completeness, recall the contrast established in Chapter~\ref{ch:nonpremixed}: a diffusion flame has no intrinsic speed and so cannot reposition itself to escape strain~\cite{lecture9}. Its incompleteness is therefore a direct consequence of the reduced residence time imposed by strain, the transport is affected immediately, and the extinction strain rates are smaller in magnitude than for premixed flames. Differential diffusion is not directly coupled to the stretch as in the premixed case, but acts through the mixing field $Z$~\cite{lecture9}.

\section{Intrinsic flame instabilities}

Even without external turbulence a premixed flame can spontaneously develop \textbf{cellular} structures, increasing its surface area and hence the fuel consumed per unit time, which can lead to self-acceleration and self-turbulisation of the front~\cite{oppong2022instability}. Three intrinsic instability mechanisms are distinguished~\cite{oppong2022instability}.

\paragraph{Hydrodynamic (Darrieus--Landau) instability.}
Discovered independently by Darrieus (1938) and Landau (1944), this instability arises purely from the \emph{thermal expansion} across the flame: the density drop deflects the streamlines through the front in such a way that any wrinkle is amplified, making a planar flame \emph{unconditionally unstable} to long-wavelength perturbations regardless of the chemistry~\cite{oppong2022instability}. Its strength grows with the expansion (density) ratio across the flame. Short-wavelength wrinkles are stabilised by the smoothing action of curvature/stretch (through the Markstein effect), so there is a cut-off scale and a corresponding \emph{critical Peclet number} (flame radius over thickness) above which an expanding flame becomes cellular~\cite{oppong2022instability}.

\paragraph{Thermo-diffusive instability.}
This mechanism is governed by the Lewis number of the deficient reactant and is the flame-front counterpart of the Markstein behaviour above~\cite{oppong2022instability}. Where the flame bulges toward the reactants (convex to the fresh gas), preferential diffusion focuses the fast-diffusing deficient reactant into the bulge: for $\Le<1$ this locally raises the temperature and burning velocity, amplifying the bulge --- the flame is thermo-diffusively \emph{unstable} and forms cells; for $\Le>1$ the opposite occurs and perturbations are damped, so the flame is \emph{stable}~\cite{oppong2022instability,lecture9}. Lean hydrogen, with its very low Lewis number, is strongly thermo-diffusively unstable, which is why lean turbulent hydrogen flames exhibit the spectacular cellular burning structures noted in Chapter~\ref{ch:governing}~\cite{day2009cellular,lecture8}.

\paragraph{Body-force (Rayleigh--Taylor) instability.}
Buoyancy acting across the density jump can also destabilise a flame, depending on orientation; it is generally secondary to the hydrodynamic and thermo-diffusive mechanisms for the flames of interest here~\cite{oppong2022instability}.

\paragraph{Consequences.}
When a flame becomes unstable and cellularised, its wrinkled area increases, so more fuel is consumed per unit time and the effective (turbulent) burning velocity rises --- the instabilities \emph{intensify} combustion~\cite{oppong2022instability}. For a spherically expanding flame this manifests as self-acceleration once the critical Peclet number is exceeded, with the flame surface becoming increasingly fractal. These effects are mild for heavy hydrocarbons (near $\Le\approx1$) but pronounced for hydrogen and syngas, and they couple with the externally imposed turbulence of Chapters~\ref{ch:turbulence} and~\ref{ch:turbcomb}; managing them is central to the design of hydrogen combustors in Chapter~\ref{ch:hydrogen}~\cite{oppong2022instability,lecture14}.
